\section{Mathematical preliminaries}\label{sec:prelim}

A small set of tools recurs in the later sections. We collect them here with the conventions of this review, so that every statement about branches, damping and dispersion rests on a single definition.

\subsection{Fourier--Laplace conventions}
Fields depend on $\exp[\ii(kx-\omega t)]$ and the Laplace variable is $s=-\ii\omega$, so that $\partial_t\to s$ and $\partial_x\to \ii k$. For a causal response, one that vanishes before the excitation, the transfer functions $G(s)$ and the resulting $k(\omega)$ are analytic in the upper half of the complex $\omega$-plane ($\operatorname{Re}s>0$). Two consequences are used below. The wave number of the uncoupled thermal branch follows from $k^2=-\rho c\,s/G(s)$. The appropriate root is the one whose spatial amplitude does not grow, $\operatorname{Im}k\ge0$, with $\operatorname{Re}k>0$ for propagation in $+x$.

\subsection{Branches and the choice of root}
For the rational heat laws treated in this review, in a linear homogeneous medium, the dispersion relation is a polynomial in $K=k^2$ with coefficients that depend on $s$; for exact delays, fractional powers or hereditary kernels those coefficients are transcendental or multivalued functions of $s$, and the rules below apply with the corresponding branch cuts. Each root $K_j(s)$ therefore gives two wave numbers $\pm\sqrt{K_j}$. Three rules are adopted: (i) for each $K_j$, take the square root with $\operatorname{Im}k\ge0$; (ii) a ``branch'' is a $k(\omega)$ followed continuously in $\omega$, because the labels ``elastic-type'' and ``thermal-type'' mean something only away from hybridisation, where roots approach each other; (iii) near hybridisation, which in the Lord--Shulman model occurs for $\tau\omega_*\approx1$, tracking by continuity on a sufficiently fine frequency grid is necessary, and selecting the root nearest an expected value fails. Both failures occurred while producing the figures of this review.

\subsection{Spatial and temporal damping}
Loss admits two measures. For real $\omega$ and complex $k$, $Q_x^{-1}=2\operatorname{Im}k/\operatorname{Re}k$; for real $k$ and complex $\omega$, $Q_t^{-1}=2|\operatorname{Im}\omega|/|\operatorname{Re}\omega|$. At weak damping the two are related by the group-to-phase velocity ratio, $Q_t^{-1}\simeq(v_g/v_p)\,Q_x^{-1}$. In the clamped-bar example of \cref{sec:bar} the classical peak $Q_t^{-1}=0.0488$ at $\varepsilon=0.1$ compares with the Debye value $\varepsilon/2=0.05$ of the spatial measure. The difference is a few percent, of the order $\varepsilon$ that the relation predicts. Resonator experiments measure $Q_t$ and wave-propagation experiments $Q_x$, and the two should not be compared without the conversion.

\subsection{Causality and the sum rule}
A causal response obeys the Kramers--Kronig relations, which tie dispersion to loss: a medium with a loss peak must show a change of speed. At weak damping the relation takes the form of a sum rule,
\begin{equation}\label{eq:sumrule}
\int_0^\infty Q_x^{-1}\,d\ln\omega=\pi\,\frac{v_0-v_\infty}{\bar v},
\end{equation}
where $v_0$ and $v_\infty$ are the phase speeds at zero and infinite frequency and $\bar v$ is their mean. For the Debye form $Q^{-1}=\varepsilon x/(1+x^2)$ the left side is $\pi\varepsilon/2$ and the right side is $\pi(\sqrt{1+\varepsilon}-1)\approx\pi\varepsilon/2$, consistent with the adiabatic and isothermal speeds of \cref{sec:classical}. We checked \cref{eq:sumrule} numerically on the exact coupled classical dispersion relation (script \path{scripts/check_kk.py}). At $\varepsilon=0.01$ the integral is $1.562994\times10^{-2}$ against $1.562991\times10^{-2}$ on the right; at $\varepsilon=0.001$ the two sides are $1.570009\times10^{-3}$ and $1.570011\times10^{-3}$. The rule is therefore a practical consistency test on any computed or measured loss curve: a model that predicts a loss peak without a corresponding speed change, or the reverse, either violates causality or has been truncated inconsistently.

\subsection{Support of the solution and finite signal speed}
The Laplace transform of a solution supported in $x\le ct$ decays as $e^{-sx/c}$ for $\operatorname{Re}s\to\infty$. The converse also holds: a transfer function with $k(s)\sim s/c$ at large $s$ implies a front travelling at $c$. That is the origin of the criterion used in \cref{sec:hyperbolic}: finite signal speed requires $sG(s)$ to tend to a finite, non-zero limit. Models with $k_w\sim\sqrt{s}$ at large $s$ (classical, DPL, TPL, GN-III, and the Guyer--Krumhansl law of \cref{sec:kinetic}) have no front, and a transient problem for them has a solution that is nonzero everywhere immediately after the excitation, however small.

\subsection{Complex roots of the delay law}
When $G$ contains an exponential, as in the exact delay law of \cref{sec:delayfid}, $k^2=-\rho c\,s/G(s)$ is an entire function times $s$ and has infinitely many branches. Truncations replace it by a rational function with finitely many roots. Roots of the delay law that the truncation does not contain, and spurious roots that it introduces, plausibly explain why truncated and exact laws differ at $\omega\tau\gtrsim1$, but this has to be checked for each truncation.

\begin{summarybox}[Summary]
The Laplace/Fourier convention ($s=-\ii\omega$, $\Imag k\ge0$ for decay) fixes branch choices once and for all. Damping, attenuation and quality factor are $Q^{-1}=2\Imag k/\Real k$ for spatial problems and $2|\Imag\omega|/\Real\omega$ for temporal problems. Causality follows from analyticity in the upper half-plane; support properties (finite front) are read from the large-$s$ behaviour of $G(s)$. The exact delay law has infinitely many branches; truncations keep finitely many and can introduce spurious roots.
\end{summarybox}
